\section{Centers and arbitrary Hilbert spaces}\label{sec:reductions}

We now pass from \Cref{thm:factors} to arbitrary von Neumann
algebras.  There are two steps.  On a separable Hilbert space, aligning
the centers permits a measurable assembly of the factor implementers.
On an arbitrary Hilbert space, separable restrictions instead supply
a completely positive map into the neighboring algebra.  The UCP
criterion then gives a unitary on the original Hilbert space.
The uniformity of \Cref{thm:factors} is essential in both steps.

\subsection{Aligning the centers}

\begin{lemma}\label{lem:centers-close}
For unital von Neumann algebras $M,N\subseteq\B(H)$,
\[
 d\bigl(Z(M),Z(N)\bigr)\leq 3d(M,N).
\]
Consequently, whenever $d(M,N)<g<1/303$, there is a unitary
$w\in\B(H)$ such that
\[
 wZ(M)w^*=Z(N),\qquad \norm{w-I}\leq450g.
\]
For $N_0=w^*Nw$ one then has $Z(N_0)=Z(M)$ and
$d(M,N_0)<901g$.
\end{lemma}

\begin{proof}
Fix $z\in Z(M)_1$ and choose $y\in N_1$ with
$\norm{z-y}<g$.  For $v\in\U(N)$, choose $a\in M_1$ with
$\norm{v-a}<g$.  Since $z$ commutes with $a$,
\[
 \norm{[z,v]}<2g,
 \qquad
 \norm{vyv^*-z}
 \leq\norm{y-z}+\norm{v z v^*-z}<3g.
\]
Dixmier's averaging theorem for von Neumann algebras
\cite[Definition~III.2.5.16 and Theorem~III.2.5.18]{BlackadarOperatorAlgebras}
says that the
norm-closed convex hull of $\{vyv^*:v\in\U(N)\}$ meets $Z(N)$.
An element of this intersection is a contraction and is within $3g$
of $z$.  Interchanging $M,N$ and then letting $g\downarrow d(M,N)$
proves the gap estimate.

The centers are amenable.  Apply the mutual conclusion of
\Cref{thm:amenable}, with near-inclusion bound $3g$, to obtain $w$.
Finally, conjugation by $w$ changes an ordinary gap by at most
$2\norm{w-I}$, giving the asserted bound for $N_0$.
\end{proof}

\subsection{The common-center decomposition}

We use the following precise form of central decomposition.  If $M$
acts on a separable Hilbert space and $Z=Z(M)$, there is a standard
measure space $(X,\mu)$ and a measurable field of separable Hilbert
spaces such that
\[
 H=\int_X^\oplus H_t\,d\mu(t),\qquad
 Z=L^\infty(X,\mu),\qquad
 M=\int_X^\oplus M_t\,d\mu(t),
\]
where $M_t$ is a factor almost everywhere.  The measure may be chosen
finite and completed.  A bounded measurable field with value in $M_t$
almost everywhere represents an element of $M$, and its norm is the
essential supremum of its fiber norms.  These are the central case of
\cite[III.1.6.2--4]{BlackadarOperatorAlgebras}.  For two algebras with
the same center, we use the same decomposition of $H$ and the same
diagonal algebra.  Countably generated $C^*$-subalgebras can be
represented simultaneously on the fibers;
see also \cite[III.5.1.10--14]{BlackadarOperatorAlgebras}.

\begin{lemma}\label{lem:fiber-gap}
Let $M,N\subseteq\B(H)$ act on a separable Hilbert space and have the
same center.  In their common central decomposition,
\[
 d(M_t,N_t)\leq d(M,N)
 \quad\text{for almost every }t.
\]
\end{lemma}

\begin{proof}
Choose separable unital $C^*$-subalgebras $A\subseteq M$ and
$B\subseteq N$ whose strong closures are $M,N$.  Such subalgebras
exist because the strong topology on each bounded operator ball is
metrizable and separable when $H$ is separable.  Choose countable
norm-dense subsets $\{a_j\}\subseteq A_1$ and
$\{b_j\}\subseteq B_1$.  Off one null set, their fiber representations
satisfy
\[
 M_t=A(t)'',\qquad N_t=B(t)''.
\]

We first check the unit-ball density needed below.  The map
$A\to A(t)$ is a quotient $*$-homomorphism, possibly with a nonzero
kernel.  A contraction in $A(t)$ has lifts of norm at most $1+s$ for
every $s>0$.  Rescaling these lifts to contractions shows that the
image of $A_1$ is norm dense in $A(t)_1$.  It follows that
$\{a_j(t)\}$ is norm dense in $A(t)_1$ and, by Kaplansky density,
strongly dense in $(M_t)_1$.  The same argument applies to $B$.

Put $d_0=d(M,N)$.  For every $j,n\geq1$, choose
$y_{j,n}\in N_1$ such that
\[
 \norm{a_j-y_{j,n}}<d_0+1/n.
\]
The essential-supremum norm formula and the countability of these
choices give a common conull set on which
\[
 y_{j,n}(t)\in(N_t)_1,
 \qquad
 \norm{a_j(t)-y_{j,n}(t)}\leq d_0+1/n
 \quad(j,n\geq1).
\]
Fix such a $t$, an $x\in(M_t)_1$, and an $n$.  Choose a net of
$a_j(t)$ converging strongly to $x$.  Along a subnet, the corresponding
$y_{j,n}(t)$ converge weakly to a contraction $y\in N_t$.
Operator-norm balls are weak-operator closed, so
$\norm{x-y}\leq d_0+1/n$.  Letting $n$ increase gives
$\dist(x,(N_t)_1)\leq d_0$.  Make the same countable choices in the
opposite direction and discard their null exceptions.  This proves
both near inclusions on one conull set.
\end{proof}

\begin{lemma}\label{lem:measurable-implementers}
In the setting of \Cref{lem:fiber-gap}, suppose that for some $a>0$
and almost every $t$ there is a unitary $v_t\in\B(H_t)$ satisfying
\[
 v_t M_t v_t^*=N_t,\qquad\norm{v_t-I}<a.
\]
Then there is a unitary $v\in\B(H)$ such that
$vMv^*=N$ and $\norm{v-I}\leq a$.
\end{lemma}

\begin{proof}
Work on the underlying standard Borel model of $X$, retaining the
countable contraction fields $a_j(t),b_j(t)$ from the preceding proof.
Choose Borel versions of the dimension function, the trivializations
on its countably many level sets, and the countably many contraction
fields.  Discard a Borel null set containing all exceptions to these
identifications, the asserted density and algebraic properties, and
the implementer hypothesis.  On each remaining dimension stratum,
the fibers are identified with one fixed separable Hilbert space $K$.
Zero-dimensional fibers contribute nothing to the direct integral
and may be omitted.  It suffices to select implementers on each of
these standard Borel strata.

The unitary group $\U(\B(K))$, with the strong-$*$ topology, is
Polish.  Fix a dense sequence $(\xi_l)$ in the unit sphere of $K$.
Consider the relation on pairs $(t,u)$ defined by
$\norm{u-I}\leq a$ and the following two families of conditions:
\begin{align*}
 &\text{for every }j,m,k\text{ there is }q\text{ such that }
 \max_{l\leq m}
 \norm{\bigl(u a_j(t)u^*-b_q(t)\bigr)\xi_l}<1/k,\\
 &\text{for every }j,m,k\text{ there is }q\text{ such that }
 \max_{l\leq m}
 \norm{\bigl(u^* b_j(t)u-a_q(t)\bigr)\xi_l}<1/k.
\end{align*}
All these predicates are Borel.  Indeed, multiplication of bounded
operator fields and strong-$*$ unitary fields is Borel, and the norm
condition on $u-I$ is a countable collection of vector-norm tests.
The strong density of the contraction fields shows that the first
family is equivalent to $uM_tu^*\subseteq N_t$, and the second to the
reverse inclusion.  Thus the relation describes exactly the
implementers with the imposed bound.  Its sections are nonempty by
hypothesis.

The Jankov--von Neumann uniformization theorem
\cite[Theorem~18.1]{Kechris1995} gives a selection measurable for the
sigma algebra generated by analytic sets, for an analytic relation
between standard Borel spaces; our relation is Borel and hence analytic.
Analytic sets are universally measurable
\cite[Theorem~21.10]{Kechris1995}, so the selector is measurable for the
completed measure $\mu$.  Its direct integral is
a unitary $v$ with $\norm{v-I}\leq a$.  For $x\in M$, the field
$v_t x_t v_t^*$ is a bounded measurable field in $N_t$, so the
field-reconstruction part of central decomposition puts $vxv^*$ in
$N$.  The reverse inclusion follows in the same way.
\end{proof}

\begin{theorem}\label{thm:separable-algebras}
For every $a>0$ there is $\delta_{\mathrm{sep}}(a)>0$ such that
whenever $M,N$ are unital von Neumann algebras on a separable Hilbert
space and $d(M,N)<\delta_{\mathrm{sep}}(a)$, there is a unitary $u$
with
\[
 uMu^*=N,\qquad \norm{u-I}<a.
\]
The tolerance is independent of the algebras and the representation.
\end{theorem}

\begin{proof}
The zero Hilbert space is immediate.  By \Cref{thm:factors}, choose
$\tau>0$ such that any pair of factors on a separable Hilbert space
at distance less than $\tau$ has an implementer within $a/2$ of the
identity.  Choose $\delta_{\mathrm{sep}}(a)>0$ so small that
\Cref{lem:centers-close} applies with $g=\delta_{\mathrm{sep}}(a)$ and
\[
 450\delta_{\mathrm{sep}}(a)<a/2,
 \qquad 901\delta_{\mathrm{sep}}(a)<\tau.
\]
For a pair satisfying the hypothesis, align the centers by the
unitary $w$ in that lemma, and set $N_0=w^*Nw$.  The common-center
pair $M,N_0$ has distance less than $\tau$.  By
\Cref{lem:fiber-gap}, almost every fiber pair has distance less than
$\tau$ as well.  Apply the factor theorem on each fiber and then
\Cref{lem:measurable-implementers}, with bound $a/2$.  The resulting
unitary $v$ satisfies
\[
 vMv^*=N_0,\qquad\norm{v-I}\leq a/2.
\]
Now $u=wv$ implements the original conjugacy and
$\norm{u-I}\leq\norm{w-I}+\norm{v-I}<a$.
Every tolerance used above is independent of the fibers and of $H$.
\end{proof}

\subsection{Separable tests in an arbitrary representation}

The preceding theorem provides one uniform unitary bound for all
separable restrictions.  These restrictions will not be chosen
compatibly.  Instead, finite collections of operator coefficients
will produce a UCP limit map, to which \Cref{prop:ucp} applies.

\begin{lemma}\label{lem:separable-tests}
Let $M,N\subseteq\B(H)$ be unital von Neumann algebras, and let
$g>d(M,N)$.  Given finite sets $E\subseteq M$, $F\subseteq N$, and
$V\subseteq H$ with $V$ containing a nonzero vector, there are
separable unital $C^*$-subalgebras $A\subseteq M$, $B\subseteq N$
and a nonzero separable closed subspace $K\subseteq H$ such that:
\begin{enumerate}[label=\textup{(\roman*)},leftmargin=*]
 \item $E\subseteq A$, $F\subseteq B$, and $V\subseteq K$;
 \item $K$ reduces both $A$ and $B$;
 \item $d(A,B)<g$ and
 \[
  d\bigl((A|_K)'',(B|_K)''\bigr)<g.
 \]
\end{enumerate}
Here the restriction homomorphisms are not required to be faithful.
\end{lemma}

\begin{proof}
Choose $r$ with $d(M,N)<r<g$, and start with
$A_0=C^*(E,I)$, $B_0=C^*(F,I)$.  Suppose that separable unital
$A_j\subseteq M$, $B_j\subseteq N$ have been constructed.  For
each element of a countable norm-dense subset of $(A_j)_1$, choose
an $N_1$-approximant within $r$, and adjoin all these approximants
to $B_j$.  In the same way, adjoin to $A_j$ approximants in $M_1$
for a countable dense subset of $(B_j)_1$.  Let $A_{j+1},B_{j+1}$
be the unital $C^*$-algebras generated by these additions.  Put
\[
 A=\overline{\bigcup_j A_j}^{\,\norm{\cdot}},
 \qquad B=\overline{\bigcup_j B_j}^{\,\norm{\cdot}}.
\]
They are separable and contain the required tests.  The stage unit
balls are norm dense in the unit balls of these closures: approximate
in norm and rescale if necessary.  The choices of approximants
therefore give $d(A,B)\leq r<g$.

Let $K=\overline{\operatorname{span}(C^*(A,B)V)}$.
Because the algebra is separable and
$V$ finite, $K$ is separable.  It contains $V$ and is reducing for
both algebras.  To transfer the gap through restriction, take
$x\in(A|_K)_1$.  For every $s>0$ there is a quotient lift
$a\in A$ with $a|_K=x$ and $\norm{a}\leq1+s$.  The contraction
$a/(1+s)$ has a $B_1$-approximant within $r$ plus arbitrary slack.
Restricting this approximant gives
\[
 \dist\bigl(x,(B|_K)_1\bigr)
 \leq \frac{s}{1+s}+r.
\]
Let $s\downarrow0$ and reverse the roles.  Thus
$d(A|_K,B|_K)\leq r$.

Finally, for a contraction in $(A|_K)''$, Kaplansky density gives
a strongly convergent net of contractions from $A|_K$.  Approximate
these in $(B|_K)_1$ with errors at most $r$ plus arbitrary slack,
and take a bounded weak-operator cluster subnet.  Its limit lies in
$((B|_K)'')_1$ and has the same error bound.  Interchanging the
algebras proves the final assertion.
\end{proof}

\begin{proof}[Proof of \Cref{thm:main}]
Let $\eps>0$.  The conclusion of \Cref{prop:ucp} is uniform in the
algebras and the Hilbert space, and its implementing-unitary bound
tends to zero with the cb error.  Choose $a>0$ and $g_0>0$ so small
that this proposition gives an implementer within $\eps$ of the
identity whenever the reverse ordinary gap is at most $g_0$ and the
cb error of the UCP map is at most $2a$.  All its smallness
requirements are included in this choice.  Choose
\[
 0<\delta<\min\{g_0,\delta_{\mathrm{sep}}(a)\}.
\]
We prove the theorem for any $M,N\subseteq\B(H)$ with
$d(M,N)<\delta$.

If $H=\{0\}$, take its identity unitary.  Otherwise fix a unit
vector $\xi_0\in H$ and a number $g$ satisfying
\[
 d(M,N)<g<\delta_{\mathrm{sep}}(a).
\]
Consider the directed set of quadruples
$i=(E_i,F_i,V_i,k_i)$, where $E_i\subseteq M$, $F_i\subseteq N$,
and $V_i\subseteq H$ are finite, $\xi_0\in V_i$, and
$k_i\in\N$ is positive.  The order is inclusion in the three finite
sets and the usual order in $k_i$.

For each $i$, apply \Cref{lem:separable-tests} to obtain
$A_i,B_i,K_i$.  Set
\[
 P_i=(A_i|_{K_i})'',\qquad Q_i=(B_i|_{K_i})''.
\]
These are unital von Neumann algebras on the nonzero separable space
$K_i$, and $d(P_i,Q_i)<g$.  By \Cref{thm:separable-algebras}, there
is a unitary $u_i\in\B(K_i)$ with
\[
 u_iP_iu_i^*=Q_i,\qquad\norm{u_i-I}<a.
\]
Extend it by the identity on $K_i^\perp$ to a unitary $U_i\in\B(H)$.
For all matrix sizes, conjugation satisfies
\begin{equation}\label{eq:finite-test-cb}
 \norm{\Ad(U_i)|_M-\id_M}_{\cb}\leq2\norm{U_i-I}<2a.
\end{equation}

We next record what the local conjugacy says about the finite tests
on the original space.  If $x\in E_i$, then $K_i$ reduces $x$, and
the restriction of $U_ixU_i^*$ belongs to $Q_i$ with norm at most
$\norm{x}$.  Kaplansky density approximates this restriction on
the finitely many coefficients from $V_i$ by an element of $B_i|_{K_i}$
of norm at most $\norm{x}$ (with arbitrarily small slack if needed).
Lift that element through the restriction quotient.  Choosing the
slack small enough gives $b_{i,x}\in B_i\subseteq N$ such that
\begin{equation}\label{eq:finite-test-lift}
 \norm{b_{i,x}}\leq\norm{x}+1,
 \qquad
 \abs{\ip{(U_ixU_i^*-b_{i,x})\xi}{\eta}}<1/k_i
 \quad(\xi,\eta\in V_i).
\end{equation}
The displayed coefficients are coefficients on $H$, since all
the vectors belong to $K_i$.  At this stage the whole conjugated
algebra need not lie in $N$; \eqref{eq:finite-test-lift} is the
property that will survive as the tests increase.

For each $x\in M$, the operators $U_ixU_i^*$ lie in the
weak-operator compact ball of radius $\norm{x}$ in $\B(H)$.
Compactness of the product of these balls gives a subnet for which
\[
 F(x)=\operatorname*{WOT\!\!-\!lim}_i U_ixU_i^*
 \quad\text{exists for every }x\in M.
\]
Linearity and unitality pass to this limit.  At each matrix size,
positive operators form a weak-operator closed cone, so $F$ is
completely positive.  Also, the norm ball of radius $2a\norm{X}$
is weak-operator closed for every fixed matrix $X$ over $M$.
Taking limits in \eqref{eq:finite-test-cb} therefore gives
\begin{equation}\label{eq:ucp-cluster-close}
 \norm{F-\id_M}_{\cb}\leq2a.
\end{equation}

It remains to verify the range of $F$.  Fix $x\in M$.  Cofinality
of the subnet implies that eventually $x\in E_i$, so choose
$b_{i,x}$ as in \eqref{eq:finite-test-lift} on this tail.  Every
fixed pair $\xi,\eta\in H$ eventually belongs to $V_i$, while
$k_i\to\infty$.  Hence \eqref{eq:finite-test-lift} and the definition
of $F$ imply $b_{i,x}\to F(x)$ in the weak-operator topology.
As $N$ is weak-operator closed, $F(x)\in N$.  This holds for every
$x\in M$, and proves that $F:M\to N$ is UCP.  Normality of this
cluster map is neither asserted nor needed.

The original inclusion $M\hookrightarrow\B(H)$ is faithful,
normal, and unital.  Its range has reverse ordinary gap less than
$\delta<g_0$ from $N$, and \eqref{eq:ucp-cluster-close} supplies the
cb bound.  Apply \Cref{prop:ucp}.  It yields a unitary $u\in\B(H)$
with $uMu^*=N$ and $\norm{u-I}<\eps$.

The choices of $a,g_0$, and $\delta$ depend only on $\eps$.
The intervening choices of finite tests, separable subspaces, and
quotient lifts change none of their bounds.  Thus the same tolerance
works for all algebras, all representations, and all Hilbert spaces,
as claimed.
\end{proof}
